\section{Determinant identities}\label{sec:preliminaries}

We begin with the two determinant formulas used throughout the paper.
The first is the finite-dimensional identity behind
\eqref{eq:intro-finite-zeta}.

Let \(A\) be a complex \(d\times d\) matrix, and set
\[
  a_n(A):=\Tr(A^n),
  \qquad
  p_n(A):=\frac{1}{n}\sum_{d\mid n}\mu(d)\,a_{n/d}(A)
  \qquad (n\ge 1).
\]

\begin{proposition}
  \label{prop:prelim-trace-primitive}
  For every complex \(d\times d\) matrix \(A\),
  \begin{equation}\label{eq:zeta-finite-matrix}
    \det(I-zA)^{-1}
    =\exp\!\left(\sum_{n\ge 1}\frac{\Tr(A^n)}{n}z^n\right)
    =\prod_{n\ge 1}(1-z^n)^{-p_n(A)}
  \end{equation}
  as an identity of holomorphic germs at \(z=0\), or equivalently as an
  identity of formal power series. If \(A\) is a primitive
  \(\{0,1\}\)-matrix, then \(p_n(A)\) counts the primitive periodic
  orbits of length \(n\) for the associated subshift
  \cite{BowenLanford1970Zeta,LindMarcus1995}.
\end{proposition}

\begin{proof}
  This is classical. By the Newton identities,
  \[
    \det(I-zA)^{-1}
    =\exp\!\left(\sum_{n\ge 1}\frac{\Tr(A^n)}{n}z^n\right).
  \]
  Since
  \(
    a_n(A)=\sum_{d\mid n}d\,p_d(A)
  \),
  we obtain
  \[
    \sum_{n\ge 1}\frac{a_n(A)}{n}z^n
    =-\sum_{m\ge 1}p_m(A)\log(1-z^m).
  \]
  Exponentiating gives \eqref{eq:zeta-finite-matrix}. The analytic
  identity holds for
  \[
    \abs{z}<\min\{1,\|A\|^{-1}\}.
  \]
  The formal version is immediate, and the final statement is standard
  in symbolic dynamics.
\end{proof}

The trace-class version is equally standard; we record it mainly to fix
notation.

\begin{proposition}
  \label{thm:fredholm-witt-bridge}
  Let \(H\) be a complex Hilbert space and let \(T\in \sone(H)\). Write
  \[
    Z_T(z):=\det(I-zT)^{-1},
    \qquad
    a_n(T):=\Tr(T^n),
  \]
  and
  \[
    p_n(T):=\frac{1}{n}\sum_{d\mid n}\mu(d)\,a_{n/d}(T)
    \qquad (n\ge 1).
  \]
  Then, for \(\abs{z}<\min(1,\|T\|^{-1})\),
  \[
    Z_T(z)
    =\exp\!\left(\sum_{n\ge 1}\frac{a_n(T)}{n}z^n\right)
    =\prod_{n\ge 1}(1-z^n)^{-p_n(T)},
  \]
  where
  \(
    (1-z^m)^{-p_m(T)}
    :=\exp\!\bigl(-p_m(T)\log(1-z^m)\bigr)
  \)
  and \(\log(1-z^m)\) is taken in the principal branch on \(\abs{z}<1\).
  In particular, the identity holds as an equality of holomorphic germs
  at \(z=0\).
\end{proposition}

\begin{proof}
  Grothendieck's trace formula gives
  \[
    Z_T(z)
    =\exp\!\left(\sum_{n\ge 1}\frac{\Tr(T^n)}{n}z^n\right);
  \]
  see \cite{Grothendieck1956} or \cite[Chapter~3]{Simon2005TraceIdeals}.
  Since \(a_n(T)=\sum_{d\mid n}d\,p_d(T)\),
  \begin{align*}
    \sum_{n\ge 1}\frac{a_n(T)}{n}z^n
      &=\sum_{n\ge 1}\sum_{d\mid n}\frac{d\,p_d(T)}{n}z^n \\
      &=\sum_{m\ge 1}\sum_{r\ge 1}\frac{p_m(T)}{r}z^{mr}
       =-\sum_{m\ge 1}p_m(T)\log(1-z^m).
  \end{align*}
  Exponentiating yields the product formula.
\end{proof}

\begin{remark}
  \label{rem:witt-audit-constraints}
  The coefficients \(p_n(T)\) may be viewed as Euler coordinates for the
  determinant, formally analogous to the big Witt coordinates; compare
  \cite{MetropolisRota1983,DressSiebeneicher1989}. In general they need
  not be integral and need not be non-negative.
\end{remark}
